\documentclass[11pt]{article}
\usepackage[a4paper,margin=25mm]{geometry}
\usepackage{amsmath,amssymb,amsthm}
\usepackage[hidelinks]{hyperref}
\emergencystretch=2em
\newtheorem{theorem}{Theorem}
\title{An irreducible chain whose empirical-measure\newline
rate function has no Hessian\\\large Disproof of Conjecture 00000002801}
\author{gaochengzhecpu}\date{}
\begin{document}\maketitle
\begin{abstract}
The deterministic alternating two-state Markov chain is
irreducible. Its actual empirical probability measures obey
a full good large deviation principle: the rate is zero
at the stationary probability vector and infinite elsewhere.
The zero is unique, but the rate is not even locally
real-valued within the interior of the probability simplex.
Its ordinary real Hessian therefore does not exist.
This disproves the Hessian assertion under the stated
irreducibility hypothesis, without any choice of spectral-gap
convention or Fisher-information normalization.
\end{abstract}

\section*{The actual irreducible chain}
Use states $0,1$ and transition matrix
\[
                   P=\begin{pmatrix}0&1\\1&0\end{pmatrix}.
\]
Every row is a probability vector. Moreover $P^2=I$:
distinct states communicate in one step and every state
returns to itself in two steps. Thus the finite chain is
irreducible in the usual positive-step reachability sense.
It is periodic, of period two. The source assumes
irreducibility and does not impose aperiodicity, laziness
or strictly positive transition entries.

Start at state zero. The complete path is
$X_k=0$ for even $k$ and $X_k=1$ for odd $k$.
Its actual law is the point mass at this path. For every
candidate path $w$ and every $n$, its cylinder probability
is exactly
\[
 \Pr(X_0=w_0,\ldots,X_n=w_n)
   =\boldsymbol1_{w_0=0}
        \prod_{k=0}^{n-1}P_{w_k w_{k+1}}.
\]
Indeed, all valid transitions force the unique alternating
path, while any invalid initial state or transition makes
both sides zero. Lean proves this identity for the actual
path probability measure, rather than only giving the
deterministic path a Markov label.

The unique stationary probability vector is
$\pi=(1/2,1/2)$. Direct multiplication gives $\pi P=\pi$;
conversely invariance forces the two masses equal and
normalization then forces each to be $1/2$.

\newpage
\section*{The complete empirical measure and its law}
Represent finite-state probability measures by their
two masses in $\mathbb R^2$. For the actual sample
$X_0,\ldots,X_n$, let
\[
 L_n(s)=\frac1{n+1}\sum_{k=0}^n\boldsymbol1_{X_k=s},
               \qquad s\in\{0,1\}.
\]
Both coordinates are nonnegative and their sum is one.
This is the entire empirical probability vector, not a
single selected statistic. The number of visits to state
one in the first $m$ observations is $\lfloor m/2\rfloor$.
In particular
\[
       |L_n(s)-1/2|\leq\frac1{n+1},\qquad L_n\to\pi.
\]
The actual distribution $\mu_n$ of $L_n$ is the pushforward
of the path law by the empirical-vector map. Since this
chain's path is deterministic, that pushforward is
$\mu_n=\delta_{L_n}$. Both its pushforward definition and
this Dirac identity are checked in Lean, as is its
probability-measure normalization.

\section*{A full good large deviation principle}
Use speed $n+1$, the actual number of observations. This
is the standard speed $m$ after writing $m=n+1$.
The LDP requires, for every open $G$ and every closed $F$,
\begin{align*}
 -\inf_{x\in G}I(x)&\leq
     \liminf_{n\to\infty}\frac1{n+1}\log\mu_n(G),\\
 \limsup_{n\to\infty}\frac1{n+1}\log\mu_n(F)&\leq
     -\inf_{x\in F}I(x).
\end{align*}
Here $\log0=-\infty$ and $\inf\varnothing=+\infty$.
These are the standard open/closed-set inequalities;
see \cite[Section 2]{varadhan}. Define
\[
 I(x)=\begin{cases}0,&x=\pi,\\+\infty,&x\ne\pi.\end{cases}
\]
We next verify these inequalities for \emph{all} open
and closed sets in the genuine real coordinate space.
Thus the rate is established by the actual distributions;
it is not merely assigned a convenient formula.

The use of the ambient $\mathbb R^2$ causes no boundary
artefact: all measures are supported on the closed
probability simplex. The bounds therefore also hold in
its relative topology, and the Hessian obstruction below
is shown along its interior affine coordinate.

\newpage
\section*{Proof of the LDP and failure of the Hessian}
If an open set $G$ contains $\pi$, convergence $L_n\to\pi$
makes $\mu_n(G)=1$ eventually. Its scaled logarithm is
eventually zero, as required since $\inf_G I=0$.
If $G$ excludes $\pi$, its required lower bound is simply
$-\infty$ and is automatic.

If a closed set $F$ excludes $\pi$, its open complement
contains $L_n$ eventually. Thus $\mu_n(F)=0$ eventually,
and the scaled logarithm is $-\infty$, equal to the
required upper bound. If $F$ contains $\pi$, its required
upper bound is zero, and every log probability is at
most zero. This proves the full LDP.

The rate is lower semicontinuous. Every finite sublevel
set is empty at negative levels and the singleton
$\{\pi\}$ at nonnegative finite levels, hence compact.
Thus the rate is good. Its unique zero is exactly the
actual unique stationary probability vector.

Now parametrize the interior probability line by
\[
                   q(t)=(1/2-t,\,1/2+t).
\]
For $|t|<1/2$ these are genuine probability vectors,
and $q(t)=\pi$ exactly when $t=0$. For every $\delta>0$
the value
\[
                 t=\tfrac12\min(\delta,1/2)
\]
satisfies $0<t<\delta$, $t<1/2$, and $I(q(t))=+\infty$.
Consequently there is no real-valued function agreeing
with this rate on any neighborhood of zero in this
interior coordinate. In particular it has no ordinary
real Hessian at its zero.

\begin{theorem}
An irreducible finite Markov chain can have a full good
empirical-measure LDP with a unique zero at which the
rate function has no ordinary real Hessian.
\end{theorem}
\begin{proof}
The alternating chain and rate constructed above satisfy
all these properties. The obstruction occurs inside the
probability simplex, before any candidate Hessian constant
or Fisher-information formula can be considered.
\end{proof}

The example refutes the source's Hessian clause. It does
not refute its unique-zero clause, which actually holds
here. Generalized extended-valued second epi-derivatives
are a different notion and are not the classical real
Hessian asserted in a Fisher-information identity.

\newpage
\section*{Formalization and reproducibility}
The actual transition kernel is a Mathlib probability
mass function: state $s$ moves to its Boolean opposite
with probability one. The complete path law is a genuine
PMF on infinite sample paths. Lean proves its exact
finite-dimensional Markov factorization both for the
outer measure and for the actual probability measure.

The empirical law is the actual pushforward of that
path law through the full occupation-vector map. The
two coordinates, probability normalization and convergence
to $\pi$ are proved from the exact occupation counts.
The actual transition matrix is irreducible, and $\pi$
is proved to be its unique normalized invariant vector.

The LDP definition in the Lean file quantifies over all
open and closed sets with the displayed standard
inequalities. Its scaled logarithm is extended-real:
probability zero gives $-\infty$, probability one gives
zero, and positive finite probabilities use the ordinary
real logarithm divided by the speed. In particular,
Lean's totalized value $\operatorname{Real.log}(0)=0$
is never used as the logarithm of a zero probability.

The general convergent-Dirac theorem proves the LDP,
then applies it to the actual empirical distributions.
Separate theorems prove lower semicontinuity, compact
finite level sets, the unique zero and absence of a
real-valued germ on the interior probability line.
No chosen numerical Hessian or spectral-gap definition
is used to create the contradiction.

Lean 4.19.0 and Mathlib revision\\
\texttt{c44e0c8ee63ca166450922a373c7409c5d26b00b}
are pinned. From \texttt{lean/}, run
\begin{verbatim}
lake exe cache get
lake build
lake env lean -DwarningAsError=true Main.lean
\end{verbatim}
The whole project is rebuilt afresh with strict checking.
Printed theorem dependencies contain only standard
foundational axioms; no placeholders, custom axioms or
external computation oracles are used. No auxiliary
numerical verifier is needed. Compile this paper with
\texttt{tectonic main.tex}.

Exact source, provenance and actual verification logs
are included. The author agent and parent performed
separate internal reviews of the proof and final
artifacts; no external independent review is claimed.

\begin{thebibliography}{9}
\bibitem{source} The Last Math Competition,
Conjecture 00000002801. Exact bilingual source included.
\bibitem{varadhan} S. R. S. Varadhan,
\emph{Large Deviations}, Spring 2010 lecture notes,
Section 2, General Principles.
\href{https://math.nyu.edu/faculty/varadhan/LDP/1-2.pdf}{Author-hosted notes}.
The reference fixes the standard LDP definition;
the example's complete proof is given above.
\end{thebibliography}
\end{document}
